\begin{textsample}{POS Dim 4 – al – Score 37.00 – al/al\_ef\_2.txt}  \label{ex:f4_pos_003}
ALAGOAS, UM \textbf{TERRITÓRIO} DE APRENDIZAGENS A área territorial de Alagoas é de 27.767,661 km2, compreendido por 102 municípios e conforme IBGE¹, a \textbf{população} totaliza 3.322.820 habitantes, com densidade demográfica de aproximadamente 119,31 hab/km2.
Com relevo formado por planície litorânea, planalto a \textbf{norte} e depressão no centro, o ponto mais elevado é a Serra de Santa Cruz, com 844 m acima do nível do \textbf{mar}.
Localizado na Região \textbf{Nordeste}, o Estado faz limites com Pernambuco, Sergipe, Bahia e o oceano Atlântico.
Em quase todo \textbf{território} alagoano, o clima é tropical, com temperaturas entre 18{\EmojiFont °}C e 26{\EmojiFont °}C.
O inverno é a estação em que ocorre maior índice de chuvas.
Os principais \textbf{rios} são : Ipanema, Moxotó, Mundaú, Paraíba do Meio e São Francisco.
O \textbf{território} é dividido geograficamente em três Mesorregiões : Mesorregiões de Alagoas Mesorregião do \textbf{Agreste} Alagoano Composta pelas Microrregiões de Arapiraca, Palmeira dos \textbf{Índios} e Traipu : Está localizada na área central do estado, está entre o \textbf{Sertão} e a Mata Atlântica.
Nos seus 24 municípios predominam culturas de feijão, fumo, amendoim, mandioca, milho, caju, algodão e cana-de-açúcar.
Neles, são encontrados minerais como : amianto, argila, calcário e ferro.
Mesorregião do Leste Alagoano Formada pelas Microrregiões do Litoral Norte Alagoano, Maceió, Mata Alagoana, Penedo, São Miguel dos Campos e Serrana dos Quilombo.
Com 52 municípios, é a Mesorregião mais habitada.
Há um intenso \textbf{turismo} por causa de suas belíssimas \textbf{praias}.
A \textbf{zona} da \textbf{mata} chama atenção, principalmente pela \textbf{flora} e \textbf{fauna} da Mata Atlântica.
Historicamente, a cana de açúcar é a cultura predominante.
Possui o Porto do Jaraguá e Aeroporto Zumbi dos Palmares.
Além do \textbf{turismo}, o comércio e a indústria fortalecem a economia da região.
Mesorregião do \textbf{Sertão} Alagoano Compreende as Microrregiões alagoanas do \textbf{Sertão} do São Francisco, Batalha, Santana do Ipanema e Serrana do \textbf{Sertão} Alagoano : É formada pela união de 26 municípios agrupados em quatro microrregiões.
O clima nessa Mesorregião é \textbf{semiárido}, com baixa umidade relativa do ar.
Vegetação de caatinga menos populosa.
A economia desta mesorregião está baseada principalmente na pecuária.
Desenvolve criação de caprinos e bovinos, além da cultura do cultivo de feijão e mandioca.
O povo de Alagoas tem características marcantes : resistência, alegria, resiliência e cordialidade.
Com origem étnico-racial na miscigenação entre os \textbf{índios} que já habitavam a região, os europeus - principalmente portugueses - e dos negros trazidos da África.
A maioria da \textbf{população} autodeclara -se \textbf{parda}, com significativa parcela de \textbf{brancos}.
A vida e experiências do povo alagoano, do Litoral, da Zona da Mata, do \textbf{Agreste} e do \textbf{Sertão}, caracterizam todo seu \textbf{território} com sua grande riqueza cultural, que se destaca no \textbf{artesanato}, \textbf{folclore}, culinária, arquitetura, religiosidade, produção literária, festas, danças, músicas, \textbf{flora}, \textbf{fauna}, lagunas, \textbf{rios}, \textbf{praias} e \textbf{mares}.
Pode -se dizer o quão são admirados o filé alagoano, a singeleza, a cerâmica e jóias \textbf{ecológicas} \textbf{indígenas} e quilombolas.
O sururu de capote, caldinho de massunim ou de feijão, fritada de siri, peixada, buchada, queijo manteiga, tapioca ou um arroz de tempero com um pirão de galinha de capoeira.
O \textbf{Coco} de roda, Reisado e Pastoril encantam toda gente, são muito apreciados de vaquejada e cavalhada. \textbf{Território} da liberdade dos povos \textbf{indígenas} e dos quilombolas, destacando -se filhos ilustres como : Zumbi dos Palmares, Graciliano Ramos, Aurélio Buarque de Holanda, Nise da Silveira, Jorge de Lima, Linda Mascarenhas, Djavan, Rainha Marta e muitos outros.
Os sotaques e expressões típicas como oxe, oxente, arretado, pocar, vôte, bulir, avia, catenga, resenha, mangar e muitas outras, definem também o \textbf{território} de Alagoas.
Os encantos naturais de Alagoas reforçam o \textbf{turismo} como um grande vetor da economia.
Destacam -se as belas \textbf{praias}, \textbf{lagoas}, lagunas e \textbf{rios}, especialmente toda a região banhada pelo Rio São Francisco.
É crescente o \textbf{turismo} \textbf{ecológico} e de aventura, com descobertas de trilhas, sítios \textbf{arqueológicos} e reservas da \textbf{mata} Atlântica.
O comércio e os serviços têm grandes destaques na economia do Estado.
Pelo interior, as grandes feiras livres trazem os produtos da agricultura familiar e da pecuária.
As culturas da cana-de-açúcar, arroz, feijão, macaxeira, inhame, milho, banana, abacaxi, \textbf{coco}, laranja, algodão e fumo são destaque.
O Estado também possui rebanhos bovinos, equinos, caprinos, bubalinos, piscicultura e ovinos.
A indústria alagoana destaca -se nos setores alimentício, açúcar, álcool, têxtil, químico, cloro químico, cimento, mineração, produção de petróleo e gás natural.
Mesmo com grande influência \textbf{branco} europeia, destaca -se que os povos tradicionais, \textbf{indígenas} e quilombolas, estão por todo Estado e garantem, principalmente, as identidades culturais do local.

% matched lemmas: agreste, arqueológico, artesanato, branco, coco, ecológico, fauna, flora, folclore, indígena, lagoa, mar, mata, nordeste, norte, pardo, população, praia, rio, semiárido, sertão, território, turismo, zona, índio
\end{textsample}
